\section{Uniform defect bounds at frequencies approaching zero}
\label{sec:near-origin-defects}
The compact bands of Section~\ref{sec:quantitative-phase-defects}
exclude the origin by a fixed amount. We now allow the physical
frequency to approach zero and the variation budget to grow. The
observable is the bounded compensation $h_R$, including its original
section discontinuity. Uniform ellipticity, rather than compact-band
phase rigidity, supplies the obstruction.

For $z\in\R^4$, $\sigma\in\R$, and a collision function $f$, put
\begin{equation}\label{eq:compensated-defect-definition}
 \mathcal D_{R,z,\sigma}f
       =f\circ T_R-e^{i(\sigma+z\cdot h_R)}f.
\end{equation}
The constant phase $\sigma$ is arbitrary. All norms in this section
are with respect to $\nu$ on the collision cylinder. Constants are
uniform in $R$; none depends on $\sigma$.

\subsection{A twisted middle block with untwisted ends}
\begin{lemma}[Damping on the diffusive collision scale]
\label{lem:diffusive-middle-damping}
There exist $r_0>0$, $A>1$, and $C<\infty$ with the following
properties. If $0<r=|z|\le r_0$, set
\[
 \delta=\tfrac14r^{1/12},\qquad m=\lceil A/r^2\rceil,\qquad
 P_{R,z}=\mathcal L_R M_{\exp(-iz\cdot h_{R,\delta})}.
\]
On every relevant local collision space,
\begin{equation}\label{eq:diffusive-middle-damping}
 \sup_{j\ge0}\|P_{R,z}^{j}\|\le C,
       \qquad |\ell(P_{R,z}^{m}\nu)|\le\tfrac18.
\end{equation}
Moreover
\begin{equation}\label{eq:middle-unsmoothing-error}
 r\|S_{m,R}(h_R-h_{R,\delta})\|_2
       \le C r^{1/24}\sqrt{1+|\log r|}.
\end{equation}
\end{lemma}
\begin{proof}
By Theorem~\ref{thm:joint-uniform-nondegeneracy},
$\Gamma_R=c_*D_R\ge c_*d_0I_4$. The covariance smoothing estimate
makes $\Gamma_{R,\delta}\ge(c_*d_0/2)I_4$ for small $r$.
The spectral expansion of Lemma~\ref{lem:smooth-collision-expansion}
applies since $r\delta^{-2}=O(r^{5/6})$. Its cubic error is
$O(r^3\delta^{-6})=O(r^{5/2})$. Reducing $r_0$ gives
\[
 \mathop{\rm Re}\log\lambda_{R,\delta}(-z)\le-c_1r^2
\]
for some uniform $c_1>0$. The projections and complementary powers
are uniformly bounded, proving the power bound. The stationary
principal amplitude is $1+O(r\delta^{-2})$. Choose $A$ large enough
that $2e^{-c_1A}<1/16$, and then reduce $r_0$ so that the
complementary term at $m$ is less than $1/16$. This proves the
second bound. Finally Proposition~\ref{prop:bv-correlation} gives
$\|S_m(h-h_\delta)\|_2^2\le Cm\delta(1+|\log\delta|)$.
Since $mr^2\le A+1$, it proves
\eqref{eq:middle-unsmoothing-error}. No multiplier bound for the
unsmoothed $h_R$ is asserted.
\end{proof}

\begin{theorem}[Quadratic coercivity for collision phases]
\label{thm:near-origin-circle-defect}
There exist $a_0,c_0,r_0>0$ such that, if $H\ge2$,
$0<r=|z|\le r_0$, and $r(1+\log H)\le a_0$, then
\begin{equation}\label{eq:near-origin-circle-defect}
 |q|=1,\quad \|q\|_{\mathrm{BV}}\le H
 \quad\Longrightarrow\quad
 \inf_{R\in I,\ \sigma\in\R}
        \|\mathcal D_{R,z,\sigma}q\|_1\ge c_0r^2.
\end{equation}
Here the infimum means the same uniform bound for every admissible
$q$ at each radius; $q$ itself may depend on $R,z,\sigma$.
\end{theorem}
\begin{proof}
Fix the radius and write $d=\|\mathcal D_{R,z,\sigma}q\|_1$.
Let $q_\epsilon=\mathsf S_\epsilon q$, where
$\epsilon=b_0/H$ and $b_0>0$ is a sufficiently small fixed number.
Then $|q_\epsilon|\le1$, $\|q-q_\epsilon\|_1\le Cb_0$, and
$\|q_\epsilon\|_{C^2}\le C\epsilon^{-2}$.
Choose integers
\[
 l=\lceil b_1+b_2\log H\rceil,\qquad N=m+2l,
\]
where $m$ is as in Lemma~\ref{lem:diffusive-middle-damping}.
The constants $b_1,b_2$ will be fixed before $H,r,q$ are chosen.

Iteration of the defect and invariance give
\begin{equation}\label{eq:near-origin-lower-pairing}
 \left|\int\overline q\,(q\circ T_R^N)
              e^{-iz\cdot S_{N,R}h_R}\dd\nu\right|\ge1-Nd.
\end{equation}
Indeed the corresponding unit-modulus integrand differs in $L^1$
from $e^{iN\sigma}$ by at most $Nd$.
Remove the phase from the first and last $l$ collisions. Since $h_R$
is bounded, this changes the pairing by at most $2lr\|h_R\|_\infty$.
Replacing its two endpoint factors by $q_\epsilon$ costs at most
$2\|q-q_\epsilon\|_1$. Smoothing only the middle collision observable
costs at most \eqref{eq:middle-unsmoothing-error}, by stationarity
and the unit supremum bounds on the endpoint factors.

The resulting pairing is exactly
\begin{equation}\label{eq:untwisted-end-word}
 \ell\left(M_{q_\epsilon}\mathcal L_R^l
                 P_{R,z}^m\mathcal L_R^l(\overline{q_\epsilon}\nu)\right).
\end{equation}
Replace both untwisted end powers by $\Pi_R=\nu\ell$.
The total error is bounded by
$C\epsilon^{-4}\vartheta^l$: both smooth endpoint factors cost
$C\epsilon^{-2}$, all other untwisted powers are bounded, and the
middle power is bounded by \eqref{eq:diffusive-middle-damping}.
The remaining expression is
\[
 \left|\int q_\epsilon\dd\nu\right|^2
                          \ell(P_{R,z}^m\nu),
\]
whose modulus is at most $1/8$.

Choose $b_0$ to make the endpoint smoothing error less than $1/16$;
then choose $b_1,b_2$ to make $C\epsilon^{-4}\vartheta^l<1/16$
for every $H\ge2$. Reducing $r_0$ makes the middle unsmoothing error
less than $1/16$. Finally choose $a_0$ so that
$r(1+\log H)\le a_0$ makes $2lr\|h_R\|_\infty<1/16$.
Thus the modulus in \eqref{eq:near-origin-lower-pairing} is at most
$1/2$. Since $N\le C/r^2$ under the same constraint, $Nd\ge1/2$
proves the assertion. Notice that the end blocks are untwisted only
inside this comparison; the original record is unchanged.
\end{proof}

\subsection{Complex vectors and their modulus}
\begin{theorem}[A uniform small-frequency function bound]
\label{thm:near-origin-vector-defect}
There exist $a_1,c_1,r_1>0$ such that, with
\[
 \Lambda(H,r)=1+\log(H/r),\qquad H\ge2,\quad 0<r\le r_1,
\]
$r\Lambda(H,r)\le a_1$ implies
\begin{equation}\label{eq:near-origin-vector-defect}
 \|f\|_2=1,\quad\|f\|_\infty+\|f\|_{\mathrm{BV}}\le H
 \quad\Longrightarrow\quad
 \|\mathcal D_{R,z,\sigma}f\|_2
                    \ge\frac{c_1r^2}{\Lambda(H,r)}
       \quad (|z|=r).
\end{equation}
The assertion is uniform in $R$ and $\sigma$, and allows zeros of $f$.
\end{theorem}
\begin{proof}
Put $e=\|\mathcal D_{R,z,\sigma}f\|_2$, $a=|f|$ and
$x=\|a-\int a\dd\nu\|_2$. The modulus estimate
\eqref{eq:approximate-modulus-variance} uses only the unit modulus of
the twist and the collision BV correlation estimate. It therefore
applies here with constants uniform in $R$:
\[
 x^2\le C H^2e^{-\gamma j}+xje.
\]
Given $\eta=b r^2$ with fixed small $b>0$, choose
$j\le C\Lambda(H,r)$ so that the first term is at most $\eta^2$.
Then $x\le\eta+je$. The circle truncation of
Lemma~\ref{lem:bv-circle-truncation} gives a phase $q$ with
\[
 \|q\|_{\mathrm{BV}}\le C(1+H),\qquad
 \|q-f\|_2\le3\sqrt2(\eta+je).
\]
Consequently
\[
 \|\mathcal D_{R,z,\sigma}q\|_1
       \le e+6\sqrt2(\eta+je).
\]
After decreasing $a_1,r_1$, the hypotheses of
Theorem~\ref{thm:near-origin-circle-defect} hold for the budget
$C(1+H)$. Choose $b$ so that the $\eta$ term is less than
$c_0r^2/2$. The remaining inequality gives
$e\ge c r^2/(1+j)\ge c_1r^2/\Lambda(H,r)$.
This proves the claim without dividing by $f$ on its zero set.
\end{proof}

\paragraph{The frequency being estimated.}
The twist in \eqref{eq:compensated-defect-definition} contains
$z\cdot h_R$, not just the physical one-collision displacement and
flight time. Its sum over a genuine return is exactly
$z\cdot(G_R-\bar G_R)$. This exact relation is used next to obtain
bounds for the actual induced record. The preceding theorems bound
defects of functions; neither is an integrated Fourier-tail estimate.
